El desvío por Santander ahorra 70\,600 COP de peaje y suma km extra. La ruta por Bogotá es más barata en dinero si el costo variable por km supera el ahorro dividido entre los km extra, $c_{km}^{*} = (P_{\mathrm{Bog}} - P_{\mathrm{San}}) / (d_{\mathrm{San}} - d_{\mathrm{Bog}})$ (en el caso central, 283.71 km extra, el ahorro de 70\,600 COP da un umbral de 248.8 COP/km). Con el costo variable vigente, 1\,157.0 COP/km (combustible con el precio de la CREG y el consumo de la UPME, insumos verificados), que supera 4.6 veces el umbral, la ruta por Bogotá es la más barata en dinero. Los otros costos por km (llantas, lubricantes, mantenimiento y conductor) no tienen fuente y se toman iguales a 0 (supuesto del equipo), de modo que el costo variable real no es menor que el calculado. No se incluyen tiempo ni riesgo.
